\documentclass[8pt]{beamer}
\usetheme{Madrid}
\usepackage[utf8]{inputenc}
\usepackage[spanish, es-tabla]{babel}
\usepackage{amsmath,amsfonts,amssymb,amsthm}
\DeclareMathOperator{\gr}{gr} % Grafo de una función
\DeclareMathOperator{\diam}{diam} % Diámetro
\DeclareMathOperator{\dist}{dist} % distancia
\DeclareUnicodeCharacter{2061}{}
\usepackage{algpseudocode}
\AtBeginSection[]
{
\begin{frame}
\begin{Large}
\tableofcontents[sectionstyle=shaded,currentsection=show]
\end{Large}
\end{frame}
}
\usepackage{mathrsfs}
\usepackage{color}
\usepackage{xcolor, colortbl}
\usepackage{array, multirow, multicol}
\usepackage{fancyhdr}
\usepackage{graphicx}
\usepackage{comment}
\usepackage{csquotes}
\usepackage[backend=biber,style=authoryear]{biblatex}
\addbibresource{referencias1.bib}
\usepackage{comment}
\usepackage{makecell}
\title[Entropía Relativa versus Volatilidad Local Estocástica: Comparación de las dos estrategias dinámicas de coberutra]{Entropía Relativa versus Volatilidad Local Estocástica: Comparación de las dos estrategias dinámicas de coberutra}
\author[Miriam Gutiérrez Serrano]{{\bf Autora:} Miriam Gutiérrez Serrano\\
{\bf Tutor acádemico:}Carlos Vázquez Cendón (UDC)\\ 
{\bf Tutor Empresa:} Daniel Arrieta Rodriguez y Cristina Arribas Francisco\\ 
\vspace{1cm} {\bf PFM Máster Matemática Industrial}}

\date{15 de Julio del 2026}

%%%%%% INICIO PRESENTACIÓN %%%%%%%%%%%%
\begin{document}

\begin{frame}
\titlepage
\end{frame}

%%%%%%%% INTRODUCCIÓN %%%%%%%%%%%
\begin{Large}
    \begin{frame}{\bf  {\LARGE {Introducción}}}
\begin{itemize}
    \item Valoración de productos derivados
    
\end{itemize}
\begin{itemize}
    \item Modelo de Black--Scholes, 1973
\end{itemize}
\begin{itemize}
    \item Black Monday, desplome bursátil del 19 de octubre de 1987
\end{itemize}
\end{frame}
\end{Large}




\begin{frame}{Índice}
    \begin{Large}
\tableofcontents
\end{Large}
\end{frame}


%%%%%%%%%%%%%%%%%%%%%%%%%%%%%%%%%%%%%%%%%%%
% CONCEPTOS PRELIMINARES Y BLACK--SCHOLES
%%%%%%%%%%%%%%%%%%%%%%%%%%%%%%%%%%%%%%%%%%%
\section{Conceptos preliminares y Black--Scholes}

% Slide 1
\begin{frame}{Definición}
\pause
  \begin{block}{Definición}
    Un proceso estocástico $\{X(t)\}_{t \geq 0}$ se dice martingala respecto de la filtración $\{\mathcal F_t\}_{t\geq 0}$ y de la medida $\mathbb P$ si es adaptado, integrable y verifica
    \[
        \mathbb{E}^\mathbb{P}[X(T) | \mathcal{F}_t] = X(t), \quad t \leq T. 
        \]
  \end{block}
  \pause
    \vspace{0.5em}
  \textbf{Primeros ejemplos:}
  \begin{itemize}
    \item Para cada $n \in \mathbb N$, $\mathbb{R}^n$ con la topología usual
    \item Cada espacio de Banach separable, con la topología inducida por su norma (Por ejemplo, los espacios $L^p$, $\ell^p,\dots$)
    \item Todo espacio metrizable compacto. 
    \item Todo espacio de Hausdorff compacto segundo contable.
  \end{itemize}
\end{frame}

% Slide 2
\begin{frame}{Teoremas fundamentales de Valoración Financiera}
  \begin{block}{Primer Teorema Fundamental.}
    Un mercado financiero está libre de arbitraje si y solo si existe una medida de probailidad $\mathbb Q$, equivalente a $\mathbb{P}$, bajo la cual los precios expresados en unidades del numerario son martingalas. Es decir, 
    \[
        \frac{S^{(m)}(t)}{N(t)} = \mathbb{E}^\mathbb{Q}\right[ \frac{S^{(m)}(t)}{N(t)} | \mathcal{F}_t \left], \quad m = 1, \dots, M, \, t \,leq T
        \]
  \end{block}
  \pause

  \begin{block}{Segundo Teorema fundamental.}
    Supongamos que el mercado está libre de arbitraje. Entonces, el mercado es completo si, y solo si, existe una única medida martingala equivalente.
  \end{block}
  \pause

  \begin{itemize}
      \item Para subespacios \textbf{cerrados}, la restricción de la métrica original garantiza la \textbf{completitud}.
      \item Para subespacios abiertos, puede construirse una métrica equivalente que garantice completitud.
      
        \vspace{0.5em}
      \textbf{Ejemplo.} El intervalo \textbf{abierto} $(0,1)$ como subespacio de $(\mathbb{R}, \tau_u)$ es un espacio polaco. No es completo con la métrica usual heredada de $\mathbb R$. Sin embargo, la métrica
    \[
    d_0(x,y) := d(x,y) + \left| \frac{1}{\dist(x,\mathbb R \setminus (0,1))} -  \frac{1}{\dist(y,\mathbb R \setminus (0,1)))}  \right|
    \]
    garantiza \textbf{completitud}.

    \item En ambos casos, subespacio de un espacio separable metrizable es \textbf{separable}.
  \end{itemize}
  


\end{frame}




% Slide 3
    \begin{frame}{Modelo de Black--Scholes}
    Consideramos un mercado formado por una cuenta bancaria libre de riesgo $B(t)$ y un activo con riesgo $S(t)$. La cuenta bancaria evoluciona según
    \[
        dB(t) = rB(t)\, dt,
        \]
    con $r$ tipo de interés libre de riesgo. 
    \pause
    El activo subyacente sigue, bajo la medida $\mathbb{Q}$,
    \[
        dS(t) = r S(t) dt + \sigma S(t) dW^\mathbb{Q} (t),
        \]
    donde $r$ es el interés libre de riesgo, $\sigma >0 $ es la volatilidad constante y $dW^\mathbb{Q}$ es un movimiento browniano bajo $\mathbb{P}$.
    \pause
    \begin{block}{\textbf{Teorema}. Ecuación de Black-Scholes.}
        Supongamos que el precio del subyacente sigue, bajo la medida neutral al riesgo, la dinámica
        \[
            dS(t) = r S(t) dt + \sigma S(t) dW^\mathbb{Q} (t).
            \]
        Entonces, el precio $V(S,t)$ de un derivado europeo con condición final $V(T,S) = g(S)$ satisface la ecuación
        \begin{equation}
            \frac{\partial V}{\partial t} + r \frac{\partial V}{\partial S} + \frac{1}{2}\sigma ^2 S^2 \frac{\partial ^2 V}{\partial S^2}-rV = 0. 
        \end{equation}
    \end{block}

\end{frame}

\begin{frame}{Modelo de Black--Scholes}
    \begin{block}{Teorema Feynman-Kac}
    Sea $X(t)$ un proceso que, bajo una medida $\mathbb Q$, satisface
    \[
        dX(t) = a(X,t)\frac{\partial V}{\partial X} + \frac{1}{2} b^2(X,t)\frac{\partial ^2 V}{\partial X^2} - r(X,t)V = 0, 
        \]
    con condición final $V(X,T)$, entonces
    \[
        V(X,T) = \mathbb{E}^\mathbb{Q}\left[ \exp \left( - \int_t^T r(X(s),s) ds\right)V(X,T) | \mathcal{F}_t \right]
        \]
    \end{block}
\end{frame}



\begin{frame}{Primeras propiedades}
 \begin{block}{Proposición}
     El producto contable de espacios polacos, ya sea finito o infinito, es polaco.
 \end{block}
 \pause
  \textbf{Ejemplo (Cubo de Hilbert).} El espacio $X = [0,1]^\mathbb N$ con la topología producto es un espacio polaco.
  \begin{itemize}
    \item En $[0,1]$, sea $d(x,y) = |x-y|, x,y \in [0,1]$ la métrica usual.
    \vspace{0.2cm}
    
    En $X$, se puede definir la métrica:
    \[ d(x, y) = \sum_{n=1}^\infty \frac{1}{2^n} |x_n -y_n|, \quad x,y \in X \]
    \item Esta métrica garantiza \textbf{completitud} y \textbf{separabilidad} en $X$.
  \end{itemize}
\end{frame}

\begin{frame}{Caracterización de subespacios polacos}
    \begin{block}{Definición}
         Sea $(X,\tau)$ un espacio topológico y sea $A\subseteq X$ no vacío. Decimos que: 
    \begin{itemize}
        \item A es un conjunto $G_\delta$ en $X$ si $\displaystyle A = \bigcap_{n \in \mathbb{N}} U_n$, donde los $U_n$ son abiertos en $X$.
        \item A es un conjunto $F_\sigma$ en $X$ si $\displaystyle A= \bigcup_{n \in \mathbb{N}} F_n$, donde los $F_n$ son cerrados en $X$. 
    \end{itemize}
    \end{block}
    \pause   
    \begin{block}{Proposición}
        Todo subconjunto cerrado de un espacio métrico es un $G_\delta$. Respectivamente, todo subconjunto abierto es un $F_\sigma$. 
    \end{block}
    \pause
      Los conjuntos $G_\delta$ no son necesariamente abiertos o cerrados. Consideremos el espacio $(\mathbb R, \tau_u)$ y $A = [0,1)$:
    \[
    A=\bigcap_{n=1}^{\infty} \left( \frac{-1}{n}, 1\right)
    \]
\end{frame}

\begin{frame}{Caracterización de subespacios polacos}
\pause
    \begin{block}{Proposición}
        Sean $X,Y$ espacios métricos homeomorfos. Entonces, si $Y$ es polaco, también lo es $X$. 
    \end{block}
    \pause
    \vspace{0.2cm}
    Si $X$ es un espacio polaco, hemos visto
    \begin{itemize}
        \item Subespacio cerrado de $X$, entonces polaco y además conjunto $G_\delta$
        \item Subespacio abierto de $X$, entonces polaco y además conjunto $G_\delta$.
    \end{itemize}
    \pause
    \begin{block}{Teorema}
    Sea $X$ un espacio polaco. Entonces, un subespacio de $X$ es polaco, si y solo si, es un conjunto $G_\delta$ en $X$.
  \end{block}
  \pause
  \vspace{0.35cm}
     \textbf{Ejemplo}. Los números irracionales constituyen un conjunto $G_\delta$ en $\mathbb R$ ($\mathbb{Q}$ es un conjunto $F_\sigma$ en $\mathbb{R}$) con respecto a la topología usual:
    \[\mathbb{R}\setminus\mathbb{Q}=\bigcap_{q \in \mathbb{Q}} \{q\}^c.
    \]
Por lo tanto, es subespacio polaco.
\end{frame}

\begin{frame}{Espacios fundamentales}
    \begin{enumerate}
        \item Espacio de Baire $\mathcal{N} = \mathbb{N}^\mathbb N$.
        \begin{itemize}
            \item Formado por sucesiones de números naturales.
            \item Base contable: $\mathcal{N}(n_1,\dots,n_k) = \{(m_i)\in \mathcal{N}: m_i = n_i, \, i =1, \dots, k \}$.
            \item Métrica completa que genera la topología producto:
            \[
            d'(\alpha, \beta) =
            \begin{cases}
            \displaystyle\frac{1}{\min\{k \in \mathbb{N} : \alpha_k \neq \beta_k\}} & \text{si } \alpha \neq \beta, \\[1.5ex]
            0 & \text{si } \alpha = \beta.
            \end{cases}
            \]
        \end{itemize}

        \only<2>{
        \textbf{Ejemplo.} Sean $       \alpha = (1,2,3,4,5,6,\dots),\quad \beta = (1,2,3,3,4,5,6,\dots)$. Se tiene que  $\displaystyle \alpha_4 =4 \neq \beta_4 = 3 \Longrightarrow$
            \[
             d'(\alpha, \beta) = \frac{1}{\min \{k \in \mathbb{N}: \alpha_k \neq \beta_k \}} = \frac{1}{4}
             \]
        }

        \only<3>{
        \item Espacio de Cantor $\{0,1\}^\mathbb N$.
        \begin{itemize}
            \item Formado por sucesiones de ceros y unos.
            \item Producto de $\{0,1\}$ con la topología discreta.
        \end{itemize}
        }
    \end{enumerate}
\end{frame}


%%%%%%%%%%%%%%%%%%%%%%%%%%%%%%%%%%%%%%%%%%%
% VOLATILIDAD LOCAL ESTOCASTICA
%%%%%%%%%%%%%%%%%%%%%%%%%%%%%%%%%%%%%%%%%%%
\section{Volatilidad Local Estocástica}

\begin{frame}{$\sigma$-álgebra de Borel y funciones medibles}
\begin{itemize}
    \item   Sean $(X_1,\mathcal{M}_1),(X_2,\mathcal{M}_2),...$ espacios medibles. El producto es el espacio $\left(\displaystyle \prod_n X_n,\prod_n \mathcal{M}_n\right)$. $\prod_n \mathcal{M}_n$ es la $\sigma-$ álgebra generada por los conjuntos de la forma 
\begin{equation*}
    A_1\times A_2 \times \dots \times A_N \times X_{N+1} \times X_{N+2} \times \dots
\end{equation*}
para algún $N \in \mathbb{N}$ y alguna elección del conjunto $A_n$ en $\mathcal{M}_n$, $n=1,...,N$. 
\pause
    \item $\mathcal{B}(X)$: $\sigma$-álgebra de Borel en el espacio topológico $(X, \tau)$, generada por su topología. 
    \pause
\end{itemize}

    \begin{block}{Proposición}
        Sean $X_1,X_2,...$ una colección finita o infinita de espacios métricos separables. Entonces, $\displaystyle \mathcal{B}\left(\prod_{n} X_n\right) = \prod_{n} \mathcal{B}(X_n)$, donde $\displaystyle \prod_n X_n$ está dotado de la topología producto. 
    \end{block}
\end{frame}

\begin{frame}{$\sigma$-álgebra de Borel y funciones medibles}
 
    \begin{block}{Definición}
            Sean $(X, \mathcal{M})$ e $(Y, \mathcal{N})$ dos espacios medibles. Una función $f:X \longrightarrow Y$ es medible con respecto a $\mathcal{M}$ y  $\mathcal{N}$ si, para cada $B \in \mathcal{N}$, el conjunto $f^{-1}(B) \in \mathcal{M}$. 
    \end{block}
    \begin{itemize}
        \item Es medible de Borel si para todo conjunto $B \in \mathcal B (Y)$ se tiene $f^{-1}(B) \in \mathcal{
B}(X)$.
    \end{itemize}
\pause
    \begin{block}{Proposición}
            Sean $(X, \mathcal{M})$ un espacio medible e $Y$ un espacio polaco. Para cada $n \in \mathbb{N}$, sea $f_n: X \longrightarrow Y$ una función medible con respecto a $\mathcal{M}$ y $\mathcal{B}(Y)$. Entonces, el conjunto 
    \begin{equation*}
        C=\left\{x\in X: \text{ existe } \lim_n f_n(x)\right\}
    \end{equation*}
    pertenece a $\mathcal{M}$. Además, la función $f:C \longrightarrow Y$ definida por 
    \begin{equation*}
        f(x) = \lim_n f_n(x)
    \end{equation*}
    es medible con respecto a $\mathcal{M}_C$ y $\mathcal{B}(Y)$. 
    \end{block}
    \begin{itemize}
        \item $\mathcal{M}_C$ denota la $\sigma$-álgebra definida por $\mathcal{M}_C := \{A \cap E: E \in \mathcal{M}\}$ 
    \end{itemize}
\end{frame}

\begin{frame}{Medidas de Borel}
\begin{itemize}
    \item Una medida definida en la $\sigma$-álgebra de Borel en un espacio topológico de Hausdorff $X$ se llama medida de Borel. 
\end{itemize}
\medskip 
\pause

\textbf{Ejemplo}.$(\mathbb R, \tau_u)$ es un espacio de Hausdorff en que todo conjunto abierto es un $F_\sigma$.
   \begin{itemize}
       \item Todo intervalo abierto es de la siguiente forma: 
       \[
        (a,b) = \bigcup_{n \in \mathbb{N}} \left[a + \frac{1}{n}, b - \frac{1}{n}\right],
   \]
   \item Todo abierto es unión contable de intervalos abiertos. 
   \end{itemize}
\pause

     \begin{block}{Teorema}
        Sea $X$ un espacio de Hausdorff donde cada conjunto abierto es un $F_\sigma$, y sea $\mu$ una medida de Borel finita en $X$. Entonces cada subconjunto de Borel en $X$ satisface que 
    \begin{equation}
    \label{eq: mu abierto}
        \mu(A) = \inf \{\mu(U): A \subseteq U, U \text{ es abierto}\}
    \end{equation}
    y
    \begin{equation}
    \label{eq: mu cerrado}
        \mu(A) = \sup \{\mu(F): F \subseteq A, F \text{ es cerrado}\}
    \end{equation}
    \end{block}
\end{frame}

\begin{frame}{Medidas de Borel}
    \begin{block}{Definición}
           Sea $(X,d)$ un espacio métrico. Un subconjunto $A \subset X$ se dice que es precompacto si, para cada número real $\epsilon >0$, existe un conjunto finito de puntos $x_1, x_2,\dots, x_n$ de $X$ tal que 
    \begin{equation*}
        A \subset \bigcup_{k=1}^nB(x_k,\epsilon)
    \end{equation*}
    \end{block}
    \begin{itemize}
        \item Un conjunto es precompacto si, y solo si, su clausura es compacta.
    \end{itemize}
    \pause
        \textbf{Ejemplo}. El conjunto $(0,1) \subset \mathbb R$ bajo la métrica usual es precompacto. 
        \pause
    \begin{block}{Teorema}
            Un espacio métrico $(X,d)$ es compacto si, y solo si, es completo y precompacto. 
    \end{block}
\end{frame}

\begin{frame}{Medidas de Borel}
 \begin{block}{Definición}
        Sea $X$ un espacio topológico de Hausdorff. Supongamos que $\mathcal{M}$ es una $\sigma$-álgebra sobre $X$ tal que $\mathcal{B}(X) \subseteq \mathcal{M}$. Una medida $\mu$ sobre $\mathcal{M}$ se dice regular si cumple las siguientes condiciones:
        \begin{enumerate}
            \item Todo subconjunto compacto $K \subseteq X$ satisface $\mu(K) < +\infty$.
            \item  Para todo conjunto $A \in \mathcal{M}$, se cumple $$\mu(A) = \inf \{\mu(U): A \subseteq U, U \text{ es abierto}\}\text{ (Regular exterior)}, \text{ y}$$.
            \item para todo conjunto abierto $U \subseteq X$, se cumple $$\mu(U) = \sup \{\mu(K): K \subseteq U, K \text{ es compacto} \} \text{ (Regular interior)}.$$
        \end{enumerate}
    \end{block}
    \begin{block}{Teorema}
        Toda medida de Borel finita en un espacio polaco es regular. 
    \end{block}
\end{frame}


%%%%%%%%%%%%%%%%%%%%%%%%%%%%%%%%%%%%%%%%%%%
% ENTROPIA RELATIVA
%%%%%%%%%%%%%%%%%%%%%%%%%%%%%%%%%%%%%%%%%%%
\section{Entropia Relativa}
\begin{frame}{Definición y caracterizaciones básicas}
    \begin{block}{Definición}
        Sea $X$ un espacio polaco. Un subconjunto $A$ de $X$ es analítico si existe un espacio polaco $Z$ y una función continua $f:Z \longrightarrow X$ tal que $f(Z)=A$. En otras palabras, $A$ es la imagen de un espacio polaco bajo una función continua.
    \end{block}
\pause
    \begin{block}{Propiedades}
         Sea $X$ un espacio polaco. Se tiene que
         \begin{itemize}
             \item si $A_1,A_2, \dots$ son subconjuntos analíticos de $X$, entonces  $\displaystyle \bigcup_k A_k$ y $\displaystyle \bigcap_k A_k$ son analíticos.
             \item Cada subconjunto $G_\delta$ de $X$ es analítico.
             \item Cada subconjunto $F_\sigma$ de $X$ es analítico.
         \end{itemize}
    \end{block}
    \pause
    \begin{itemize}
        \item En particular, abiertos y cerrados de $X$ son analíticos.
    \end{itemize}
\end{frame}

\begin{frame}{Relación con el espacio de Baire}
    \begin{block}{Teorema de la intersección de Cantor}
        Sea $(X,d)$ un espacio métrico completo y sea $(F_n)$  una sucesión decreciente (con respecto a la inclusión) de conjuntos cerrados no vacíos de $X$ tal que $\displaystyle \lim_{n \to \infty} \diam (F_n) = 0 $. Entonces, $\displaystyle \bigcap_{n=1}^\infty F_n$ contiene exactamente un solo punto, es decir, 
    \begin{equation*}
        \bigcap_{n=1}^\infty F_n = \{x\} \text{ para algún } x \in X
    \end{equation*}
    \end{block}
    \begin{block}{Teorema}
        Todo espacio polaco no vacío es la imagen de $\mathcal{N}$ por una función continua.
    \end{block}
\end{frame}

\begin{frame}{Relación con el espacio de Baire}

    \begin{block}{Corolario}
        Todo subconjunto analítico de un espacio polaco no vacío es la imagen de $\mathcal{N}$ por una función continua.
    \end{block}
    \pause 
    \begin{block}{Proposición}
        Sea $X$ un espacio polaco. Un subconjunto $A$ de $X$ es analítico si, y solo si, existe un subconjunto cerrado de $\mathcal{N} \times X$ cuya proyección sobre $X$ sea $A$.
    \end{block}
\end{frame}


\begin{frame}{Relación con el espacio de Cantor}
    \begin{block}{Definición}
      Un espacio topológico $(X,\tau)$ se dice de   dimensión cero si admite una base formada por conjuntos que son tanto abiertos como cerrados.  
    \end{block}
    \pause
    \vspace{0.35cm}
    \textbf{Ejemplos}.
    \begin{itemize}
        \item  Consideramos $\mathbb{Q}$ dotado de la topología heredada de $\mathbb{R}$. Los conjuntos de la forma
    \[
    (a,b) \cap \mathbb{Q} = [a,b] \cap \mathbb{Q}, \quad a,b \in \mathbb{R} \setminus\mathbb{Q},\ a < b,
    \]
    constituyen una base para la topologia.
    \item El espacio de Baire $\mathcal{N}$.
    \item El espacio de Cantor $\{0,1\}^\mathbb{N}$.
    \end{itemize}
   
\end{frame}
\begin{frame}{Relación con el espacio de Cantor}
    \begin{block}{Proposición}
    Sea $X$ un espacio separable y metrizable. Entonces es homeomorfo a un subespacio de $[0,1]^\mathbb{N}$. 
    \end{block}
    \begin{block}{Proposición}
        Sea $X$ un espacio polaco. Entonces es homeomorfo a un subconjunto $G_\delta$ de $[0,1]^\mathbb{N}$.
    \end{block}
    \pause
    \begin{itemize}
        \item O sea: $[0,1]^\mathbb{N}$ es polaco, y además, ''contiene'' a todo espacio polaco.
    \end{itemize}
    \pause
    \begin{block}{Proposición}
        Todo espacio polaco es la imagen de un espacio polaco de dimensión cero por una función inyectiva continua.
    \end{block}
    \begin{block}{Proposición}
        Todo subconjunto de Borel de un espacio polaco es la imagen de un espacio polaco de dimensión cero por una función inyectiva continua.
    \end{block}
\end{frame}

\begin{frame}{Relación entre representaciones de espacios polacos}
    Sea $X$ un espacio polaco no vacío. Hemos visto
    \begin{enumerate}
        \item Existe una función continua $f: \mathcal{N}\longrightarrow X$ sobreyectiva:
        \[
        X = f(\mathcal{N}), \quad \text{con } \mathcal{N} = \mathbb{N}^\mathbb{N}
        \]
        \item $X$ es \textbf{homeomorfo} a un subconjunto $G_\delta$ de $[0,1]^\mathbb{N}$ (espacio polaco).
        \item $X$ es la imagen de un espacio polaco de dimensión cero por una función continua e inyectiva
    \end{enumerate}
    \vspace{0.5cm}   
    
   \textbf{Resumen:}
\[
\mathcal{N} \xrightarrow{\text{cont., sobrey.}} X 
\xrightarrow[\text{homeo}]{\phantom{\text{homeo}}} 
G_\delta \subseteq [0,1]^\mathbb{N},
\quad 
Y \xrightarrow{\text{cont., biyectiva}} X
\]
\begin{itemize}
  \item \( \mathcal{N} \): espacio de Baire.
  \item \( Y \): espacio polaco de dimensión cero.
\end{itemize}

\end{frame}

\begin{comment}
    \begin{frame}{Relación con el espacio de Cantor}
    \begin{block}{Definición}
       Sea $X$ un espacio métrico. Decimos que $p$ es un punto de condensación si cada entorno de $p$ contiene una cantidad no contable de puntos. 
    \end{block}
    \vspace{0.2cm}
    \textbf{Ejemplo.} Sea $X = \mathbb R, \,\mathbb{R}\setminus \mathbb{
Q}$ con la métrica usual. Tomamos un punto $p \in X$. En efecto, es punto de condensación:
    \[
    \text{para cualquier } \epsilon>0, B(p,\epsilon) \cap X \text{ es no contable}
    \]
    \pause
    \begin{block}{Proposición}
       Sea $X$ un espacio topológico segundo contable. Si $\mathcal{V}$ es una familia de subconjuntos abiertos de $X$, entonces existe $\mathcal{V}_0 \subset \mathcal{V}$ contable tal que 
    \[
    \bigcup_{V \in \mathcal{V}} V = \bigcup_{V \in \mathcal{V}_0} V 
    \]
    \end{block}
\end{frame}
\end{comment}


\begin{frame}{Relación con el espacio de Cantor}
    \begin{block}{Teorema}
        Sea $X$ un espacio polaco, y sea $B$ un subconjunto no contable de $X$. Entonces, existe una función continua inyectiva $f:\mathcal{N}\longrightarrow X$ tal que $f(\mathcal{N}) \subseteq B$ y $B\setminus f(\mathcal{N})$ es contable.
    \end{block}
    \vspace{0.35cm}
    \pause
    Sean ahora $X$ un espacio compacto e $Y$ un espacio de Hausdorff.
    Supongamos que $f:X \longrightarrow Y$ es una función continua y biyectiva. 
     \begin{itemize}
         \item Como $X$ es compacto y $f$ es continua, $f(X)$ es compacto.
         \item $f$ es un homeomorfismo entre $X$ e $Y$.
     \end{itemize} 
    \pause
    \vspace{0.35cm}
    \begin{block}{Corolario}
       Todo subconjunto de Borel no contable de un espacio polaco contiene un subconjunto que es homeomorfo a $\{0,1\}^\mathbb{N}$ 
    \end{block}
\end{frame}

%%%%%%%%%%%%%%%%%%%%%%%%%%%%%%%%%%%%%%
% IMPLEMENTACIÓN NUMERICA
%%%%%%%%%%%%%%%%%%%%%%%%%%%%%%%%%%%%%%%

\section{Implementación numérica}
\begin{frame}{Definición y caracterizaciones básicas}
    \begin{block}{Definición}
        Sea $X$ un espacio polaco. Un subconjunto $A$ de $X$ es analítico si existe un espacio polaco $Z$ y una función continua $f:Z \longrightarrow X$ tal que $f(Z)=A$. En otras palabras, $A$ es la imagen de un espacio polaco bajo una función continua.
    \end{block}
\pause
    \begin{block}{Propiedades}
         Sea $X$ un espacio polaco. Se tiene que
         \begin{itemize}
             \item si $A_1,A_2, \dots$ son subconjuntos analíticos de $X$, entonces  $\displaystyle \bigcup_k A_k$ y $\displaystyle \bigcap_k A_k$ son analíticos.
             \item Cada subconjunto $G_\delta$ de $X$ es analítico.
             \item Cada subconjunto $F_\sigma$ de $X$ es analítico.
         \end{itemize}
    \end{block}
    \pause
    \begin{itemize}
        \item En particular, abiertos y cerrados de $X$ son analíticos.
    \end{itemize}
\end{frame}

%%%%%%%%%%%%%%%%%%%%%%%%%%%%%%%%%%%%%%%%%%%%%%%%
% CONCLUSIONES Y LÍNEAS FUTURAS DE INVESTIGACIÓN
%%%%%%%%%%%%%%%%%%%%%%%%%%%%%%%%%%%%%%%%%%%%%%%%
\section{Conclusiones}


%%%%%%%%%%%%%%%%%%%%% REFERENCIAS %%%%%%%%%%%%%%%%%%%%%%%
\begin{frame}{\bf Referencias}
\begin{thebibliography}{50}

  \bibitem{lib1:libro} Cohn, D.L.
  {\em Measure Theory: Second Edition}. Birkhäuser, New York, 2013.

  \bibitem{lib2:libro} Willard, S.
  {\em General Topology}. Addison-Wesley, 1970.
  \bibitem{lib3:libro} Munkres, J.R.
  {\em Topology. Second edition}. Pearson, 2002.

  \bibitem{araujo:apuntes} Araujo Gómez, J.
  {\em Measure Theory, Parts 1--3}. Apuntes del curso 2023/2024, Universidad de Cantabria.



\end{thebibliography}
\end{frame}



%%%%%%%%%%%%%%%%% FIN %%%%%%%%%%%%%%%%%%%%%%%
\begin{frame}
\titlepage
\end{frame}
\end{document}
